% \iffalse meta-comment
%
% Copyright (C) 2021 by Robert J Lee
%
% This file may be distributed and/or modified under the
% conditions of the LaTeX Project Public License, either
% version 1.3 of this license or (at your option) any later
% version. The latest version of this license is in:
%
% http://www.latex-project.org/lppl.txt
%
% and version 1.3 or later is part of all distributions of
% LaTeX version 2005/12/01 or later.
%
% \fi
%
% \iffalse
% <package>\NeedsTeXFormat{LaTeX2e}[2005/12/01]
% <package>\ProvidesPackage{gamebooklib}
% <package>   [2022/05/26 v1.3 gamebooklib typesetting]
%
%<*driver>
\documentclass{ltxdoc}
\usepackage{makeidx}
\usepackage[seed=1234]{lcg}
\usepackage{gamebooklib}
\usepackage{indentfirst}
\EnableCrossrefs
\CodelineIndex
\RecordChanges
\makeindex
\begin{document}
  \DocInput{gamebooklib.dtx}
\end{document}
%</driver>
% \fi
%
% \CharacterTable
%  {Upper-case    \A\B\C\D\E\F\G\H\I\J\K\L\M\N\O\P\Q\R\S\T\U\V\W\X\Y\Z
%   Lower-case    \a\b\c\d\e\f\g\h\i\j\k\l\m\n\o\p\q\r\s\t\u\v\w\x\y\z
%   Digits        \0\1\2\3\4\5\6\7\8\9
%   Exclamation   \!     Double quote  \"     Hash (number) \#
%   Dollar        \$     Percent       \%     Ampersand     \&
%   Acute accent  \'     Left paren    \(     Right paren   \)
%   Asterisk      \*     Plus          \+     Comma         \,
%   Minus         \-     Point         \.     Solidus       \/
%   Colon         \:     Semicolon     \;     Less than     \<
%   Equals        \=     Greater than  \>     Question mark \?
%   Commercial at \@     Left bracket  \[     Backslash     \\
%   Right bracket \]     Circumflex    \^     Underscore    \_
%   Grave accent  \`     Left brace    \{     Vertical bar  \|
%   Right brace   \}     Tilde         \~}
%
% \CheckSum{708}
%
% \changes{v1.0}{2021/08/10}{Initial Release}
% \changes{v1.1}{2021/08/24}{Bugfix: edge case with clashing fixed-index entries}
% \changes{v1.2}{2021/09/11}{Bugfix: footnotes set justified on last line}
%
% \GetFileInfo{gamebooklib.sty}
%
% \DoNotIndex{\let, \edef, \xdef, \relax, \newcommand}
% \DoNotIndex{\csname, \endcsname, \@M, \@MM, \@gobble}
% \DoNotIndex{\@ifnextchar, \@tmp, \@undefined, \\}
% \DoNotIndex{\arabic, \begin, \begingroup, \botmark, \break}
% \DoNotIndex{\@tmp, \@tmp@, \mytmp@, \tmp@, \tmp@@}
% \DoNotIndex{\centerline, \cs, \def, \detokenize, \else}
% \DoNotIndex{\endgroup, \equal, \expandafter, \fi}
% \DoNotIndex{\gdef, \global, \head, \hspace, \Huge, \if}
% \DoNotIndex{\ifcsname, \ifhmode, \ifnum, \ifthenelse, \ifvmode}
% \DoNotIndex{\numexpr}
% \DoNotIndex{\interlinepenalty, \interlinefootnotepenalty}
% \DoNotIndex{\large, \leftskip, \mark, \NeedsTeXFormat}
% \DoNotIndex{\newcounter, \newenvironment, \newtoks, \noexpand}
% \DoNotIndex{\noindexnt, \nopagebreak, \not, \PackageInfo, \pagebreak}
% \DoNotIndex{\par, \PassOptionsToClass, \penalty, \ProcessOptions}
% \DoNotIndex{\protected@xdef, \ProvidesPackage, \rand, \requirecommand}
% \DoNotIndex{\RequirePackage, \rightskip, \rule, \setbox}
% \DoNotIndex{\setcounter, \stepcounter, \textbf, \value}
% \DoNotIndex{\vbox, \unvbox, \vfill, \vspace, \vsize}
% \DoNotIndex{\WarningFilter, \whiledo, \Collect@Body, \CurrentOption}
% \DoNotIndex{\DeclareOption, \g@addto@macro, \the, \noindent}
% \DoNotIndex{\interfootnotelinepenalty, \marginpar}
% \DoNotIndex{\addtocounter, \refstepcounter, \label}
% \DoNotIndex{\iffalse, \iftrue, \outputpenalty, \renewcommand}
% \DoNotIndex{\noshuffle, \verbose, \endpage, \footnote}
% \DoNotIndex{\if\dots}
%
% \title{The Gamebooklib Package\thanks{This document corresponds to \textsf{Gamebooklib}~\fileversion, dated \filedate.}}
% \author{Robert J Lee \\ latex@rjlee.homelinux.org}
%
% \maketitle
%
% \begin{abstract}
%
% This package provides macros and environments to allow the user to
% typeset a series of cross-referenced, numbered ``entries'',
% shuffled into random order
%
% \end{abstract}
%
% \section{Introduction}
%
% This package was written to allow the typesetting of gamebooks.
%
% A gamebook is a book divided into ``entries'' (which may be a single
% paragraph or longer), each with a sequentially numbered value. At
% the end of each entry, a link to one or more other numbered
% entries is given; the reader selects one and they follow through
% the story in this order. Gamebooks traditionally start at the
% entry numbered~``1'' and some gamebooks have multiple endings,
% with the final section representing a success or victory for the
% reader.
%
% In particular, this package handles two technical challenges which
% are tricky to solve with \LaTeX\ directly: writing the ``entries''
% in order but presenting them in a randomised order; and presenting
% footnotes at the end of each ``entry'', numbered in the order in
% which they appear, not the order in which they are written.
%
% Aside from shuffling the paragraphs and fixing the handling of
% footnotes, this package provides little help for the actual
% typesetting of gamebook entries. For this, please consider using
% Andr\'e Miede's |gamebook| package, available on
% CTAN\footnote{\texttt{https://www.ctan.org/pkg/gamebook}}. That
% package is orthogonal to this one; both may be used together.
%
% \section*{Terminology}
%
% For the purposes of this document, the term \textit{entry} means a
% numbered section of text. This term helps describe the
% common format of a gamebook, while avoiding confusion with the terms
% \textit{paragraph} and a \textit{section}, both of which
% already have a clear definition.
%
% The term \textit{gamebook} means a book consisting of numbered
% sections, typically presented in a non-linear order with a
% tree or graph of possible reading options.
%
% \section*{Usage}
%
% To load the class, use |\usepackage[|\textit{options}|]{gamebook}|
%
% The class options are described below; |[footnotes,jukebox]| is
% normally a good choice, although there are some limitations.
%
% The user simply writes entries like this inside the document:
% \begin{verbatim}
% \begin{gentry}{codea}
% This is the first entry
%
% Turn to \turnto{codeb}
% \end{gentry}
%
% \begin{gentry}[123]{codeb}[title]
% This is the second and final entry.\par
% To play again, turn to \turnto{codea}
% \end{gentry}
%
% % Output:
% \thegentries
% \end{verbatim}
%
% Here, two entries are defined, and then output.
%
% The first entry is given a label of |gentry:codea|, and the
% \cs{turnto\{codea\}} in the second entry generates a link back to
% that label.
%
% As |entry| is a common term likely to conflict with other
% environments, the environment to declare an entry has been named
% |gentry| (gamebook entry) instead.
%
% Sometimes it is convenient, when writing an entry, to give it a fixed
% index in the text. The first optional argument in the second entry
% --- |[123]| --- fixes the entry at that position during
% shuffling. This is implemented by iterating over each entry and
% swapping it with that index. Duplicates are not currently checked
% for, but they would not be lost; if two entries are defined with the
% same fixed index, one will end up somewhere else in the text. 
%
% The first and last entries are automatically fixed without
% needing to specify the optional number.
%
% Specifying the optional number as a value greater than or equal to
% the number of entries, or less than~2, is not recommended. (Yes,
% this example used |123|, which is beyond the number of entries. This
% was done to show the syntax).
%
% The second optional argument is a title to be displayed alongside
% the numerical value.
%
% Finally, the command \cs{thegentries} will output all entries defined so
% far, handling shuffling and footnotes.
%
% \DescribeMacro{\gentryheader}\marg{counterIdx}\marg{fixedIdx}\marg{code}\marg{title}
%
% Users may also want to redefine \cs{gentryheader} to change how the
% entries are displayed. The initial version is quite basic and
% intended for debugging; users should use \cs{renewcommand} to change
% it to something more pleasant.
%
% Here's one suggestion, using packages from CTAN. It works well if you have either a
% title, or at least 2 lines of text before the first paragraph break
% inside the |gentry|.
%
%
% It also requires near the start of the document:
% \begin{verbatim}
% \usepackage{lettrine}
% \usepackage{etoolbox}
% \usepackage{color}
% \end{verbatim}
% And this anywhere before \cs{thegentries}:
% \begin{verbatim}
% \renewcommand{\LettrineFontHook}{%
%   \color[gray]{0.5}\fontfamily{ppl}\fontseries{bx}}
% \renewcommand{\gentryheader}[4]{
%  \lettrine[lines=3,lhang=0.2,loversize=0.2]{\raisebox{0.1em}{#1}}%
%      {\textbf{\Large\textbf{\raisebox{0.3em}{~#4}
%      }}}%
%      \ifstrempty{#1}{}{\linebreak\mbox{}~}}
% \end{verbatim}
% On the last line, \cs{mbox} stops the \cs{parindent} space being
% ignored; the $\sim$ brings it back up to the expected spacing.
%
% \section*{Package options}
%
% These options can be specified as a comma-separated list to the
% \cs{usepackage} line.
%
% \begin{macro}{footnote}
% \changes{v1.0}{2021/08/10}{Footnote option}
% The |footnote| option causes \LaTeX\ to output footnotes at the end
% of each entry, or the end of each page, whichever comes first after
% the footnote mark.
%
% There are some limitations, described below.
%
% Support for other packages that affect footnotes is limited. If
% using |hyperref|, it is recommended to pass |[hyperfootnotes=false]|
% to avoid broken links.
% \end{macro}
%
%
% \begin{macro}{jukebox}
% \changes{v1.3}{2022/05/26}{Feature: "jukebox" shuffle}
% Use the Jukebox Index shuffling
% algorithm\footnote{\texttt{https://github.com/robertjlee/jukeboxshuffle}}.
% This is slightly slower, but tends to reduce the number of times you
% get a  ``turn to~'' instrution referencing the next (or previous)
% entry in the original order, by modifying the shuffle to ensure that
% adjacent gentries in the input have much less chance of being
% adjacent in the final document. |jukebox| requires a \LaTeXe that
% supports \cs{numexpr}, and a minimum of 6 gentries. If fewer than
% 6~|gentry| environments are supplied before \cs{thegentries}, this
% option will only log a warning in verbose mode.
%
% The |jukebox| option guarantees no more adjacent entries than
% without the option, for a given |seed| value; it may not eliminate
% them completely unless the number of entries is large. The
% performance is linear to the number of entries.
% \end{macro}
%
% \begin{macro}{noshuffle}
% \changes{v1.0}{2021/08/10}{Noshuffle option}
% In general, it's easier to write gamebooks in a more linear fashion,
% in which related entries are kept together. But this is much less
% fun to play, as it's too easy for the reader to simply read the
% adjacent eentries to decide what to do.
% For this reason, this package shuffles the output entries by
% default. The first and last entries are never shuffled.
%
% Sometimes, perhaps for proofreading, you may not want the entries
% to be shuffled. In this case, you can use the |noshuffle| option.
% \end{macro}
%
% \begin{macro}{verbose}
% \changes{v1.0}{2021/08/10}{Verbose option}
% This macro causes the package to output information messages about
% what it's doing to the log file. This is not generally too useful, but it does
% include a mapping of the original paragraph indexes to their sorted
% positons, which may be useful to keep for handling proofreading
% corrections.
% \end{macro}
%
% \begin{macro}{endpage}
% \changes{v1.0}{2021/08/10}{Endpage option}
% Puts the last entry on a page on its own. This typically produces a
% better ``winning'' feel for reaching the last entry, but it can also
% produce documents with ugly spacing, so it's recommended to try it
% each way and see which works better for your text.
% \end{macro}
%
% \begin{macro}{seed}
% It's suggested that users specify a value for ``seed'' for stable
% builds, by adding the following before including this package:
% |\usepackage[seed=123]{lcg}|
% \end{macro}
% \section*{Footnotes}
%
% When typesetting a gamebook with footnotes, it is confusing if they
% migrate to the bottom of each page, as the footnote becomes visually
% detached from the entry to which it relates. Some effort has been taken to
% ensure that footnotes can be typeset at the bottom of the page on which
% the mark appears, or the bottom of each entry, whichever comes
% first.
%
% However, this implementation has some limitations:
%
% \begin{itemize}
% \item Footnotes are not expanded until they are typeset.
% \item As a consequence, attaching one footnote to another with the
% use of \cs{footnotemark} within a \cs{footnote} argument will not
% advance the counter in time, so \cs{footnotetext} may not behave as
% you expect unless it is also set inside the first footnote's text.
% \item Footnotes at the end of the page will not be broken across
% pages. Putting sufficient text into footnotes will cause the page
% to overflow.
% \item Where the footnote \textbf{mark} appears at the very end of a page,
% the footnote \textbf{text} may be set at the top of the subsequent
% page. \TeX\ is asked not to break the page there, but this influence
% is limited. It may be possible for the author to avoid this, eg by
% adding a rubber length to the interline spacing (it may be easier to
% simply choose a different |seed| value for the |lcg| package to
% reshuffle).
% \item If you have a lot of rubber space in the text, or
% variably-sized items, \LaTeX\ may expand the footnote at the end of
% the entry before it works out where to put the page break. In this
% case, the footnote will appear on the wrong page. The macro
% \cs{noentryfoot} is provided to allow the author to fix this case.
% \item Support for other footnotes packages may be limited and the
% behaviour of packages like |footnote|, |endnotes|, |footmisx|,
% |fnote|, |dblfnote| \textit{etc} may be affected.
% \item If using the |hyperref| package, it is recommended to pass the
% option |hyperfootnotes=false| to that package, as the footnote links
% will be incorrect.
% \item The package uses \TeX\ \cs{mark}s to indicate which footnotes
% should appear on each page. Because \TeX\ supports only one set of
% marks, this would break any other package or usage of \cs{mark}
% while a gamebook was being output.
% \item Because \cs{mark} does not escape a floating environment, such
% as |minipage|, it is likely that this footnote implementation will
% not work as expected if the gamebook or \cs{footnote} is set in a |minipage| or
% similar.
% \item Each footnote is evaluated in a group. Just conceivably, this
% might affect the behaviour of some macros that affect the document
% beyond the footnote.
% \item We rely on some non-``public'' macros, such as \cs{@mpfn} and
% \cs{@footnotetext} defined by \LaTeX\ standard document
% classes. Other document classes may override these, which would
% override this package's footnotes too.
% \end{itemize}
%
% For this reason, the improved footnote handling is only enabled if
% you pass the |footnote| option, and only while the environment is
% active.
%
% \section*{Implementation}
%\StopEventually{\PrintChanges \pagebreak[4] \PrintIndex}
%
%    \begin{macrocode}
\ProvidesPackage{gamebooklib}[2022/05/25 Gamebook by R Lee latex@rjlee.homelinux.org]
%    \end{macrocode}
% We need \LaTeXe, for the extra token registers.
%    \begin{macrocode}
\NeedsTeXFormat{LaTeX2e}[1994/06/01]
%    \end{macrocode}
% \begin{macro}{verbose}
% The package option |verbose| enables detailed logging.
% Logging is via the macro \cs{gamebook@info}, which throws away detail messages
% unless the |verbose| option is given.
%    \begin{macrocode}
\newcommand{\gamebook@info}[1]{}%
\DeclareOption{verbose}{%
  \renewcommand{\gamebook@info}[1]{\PackageInfo{gamebooklib}{#1}}%
  \gamebook@info{Gamebook Library package is logging}}%
%    \end{macrocode}
% \end{macro}
%
% \begin{macro}{endpage}
% The |endpage| option puts the last entry on its own page. This can
% work better when the last entry is about a page long, and also
% the final ``winning'' entry of the gamebook.
%    \begin{macrocode}
\newcommand\gamebook@beforelast{}
\DeclareOption{endpage}{%
  \renewcommand\gamebook@beforelast{\eject}%
}%
%    \end{macrocode}
% \end{macro}
%
% \begin{macro}{jukebox}
% The |jukebox| option defines the \cs{gamebox@jukebox} macro;
% while the macro does nothing, \cs{ifcsname} can then be used to
% determine if the option was set.
%    \begin{macrocode}
\DeclareOption{jukebox}{%
  \newcommand\gamebook@jukebox{}%
  \gamebook@info{Gamebook Library to perform jukebox index reshuffle
    pass}%
}%
%    \end{macrocode}
% \end{macro}
%
% \begin{macro}{footnote}
% The |footnote| option enables our footnote processing, to throw out
% footnotes at the end of each |gentry| so that they don't appear to be
% against subsequent entries for that page.
%
% This is genearlly recommended, unless you have a reason to turn it
% off (such as a conflicting package). It's disabled by default
% because it could cause unexpected faults.
%    \begin{macrocode}
\def\if@gamebook@footnotes{\iffalse}
\DeclareOption{footnote}{%
  \gdef\if@gamebook@footnotes{\iftrue}%
  \gamebook@info{Gamebook Library footnotes per gamebook entry}%
}%
%    \end{macrocode}
% \end{macro}
%
% \begin{macro}{noshuffle}
%    \begin{macrocode}
\def\if@gamebook@shuffle{\iftrue}
\DeclareOption{noshuffle}{%
  \gdef\if@gamebook@shuffle{\iffalse}%
  \gamebook@info{Gamebook Library entries output in order}%
}%
%    \end{macrocode}
% \end{macro}
%
% \begin{macro}{seed}
% All unknown options are passed to |lcg|, as it's our only dependency with options.
%    \begin{macrocode}
\DeclareOption*{%
  \PassOptionsToClass{\CurrentOption}{lcg}%
}%
\ProcessOptions\relax%
%    \end{macrocode}
% \end{macro}% We need to capture environment contents
%    \begin{macrocode}
\RequirePackage{environ}%
%    \end{macrocode}
% |Macroswap|: used to swap the commands that evaluate the macros
%    \begin{macrocode}
\RequirePackage{macroswap}%
%    \end{macrocode}
% |Ifthen| makes branching and loops a little easier
%    \begin{macrocode}
\RequirePackage{ifthen}
%    \end{macrocode}
% |LCG|: random numbers for the shuffle
%    \begin{macrocode}
\RequirePackage{lcg}%
%    \end{macrocode}
% |Silence| is used to suppress a warning from~|lcg| that gamebook is
% \textbf{not} wasting counter registers. |debrief| ensures that the
% user is at least told than warnings were suppressed.
%
% To see the warnings, put the following line before
% |\usepackage{gamebook}|:
%
% |\usepackage[debrief,showwarnings]{silence}|:
%    \begin{macrocode}
\RequirePackage[debrief]{silence}%
%    \end{macrocode}
%
% Let's count the entries we're reading in so we can build up token
% registers.
%    \begin{macrocode}
\newcounter{gentryctr}%
\setcounter{gentryctr}{0}%
%    \end{macrocode}
%
% \begin{environment}{gentry}\oarg{fixedIdx}\marg{code}\oarg{title}
% First, we need to parse the optional arguments.
%
% When we say \cs{begin}\{gentry\}, \LaTeX\ does some stuff ending in \cs{gentry}
% so we redefine \cs{gentry} to take an optional argument and delegate to \cs{@gentry}
%    \begin{macrocode}
\newcommand{\gentry}[1][]{\@gentry{#1}}%
%    \end{macrocode}
% \cs{gentry} just makes the first argument mandatory.
%
% NB: If you define two entries with the same code, \LaTeX\ will print
% out a "multiply defined" label warning.
%    \begin{macrocode}
\newenvironment{@gentry}[2]{%
  \xdef\gentryidx{#1}%
  \xdef\gentrycode{#2}%
  \@@gentry%
}{\ignorespacesandallpars%
  \global\let\gentryidx\@undefined%
  \global\let\gentrycode\@undefined%
  \global\let\gentryidxu\@undefined%
  \global\let\gentryidxs\@undefined%
}%
%    \end{macrocode}
% \begin{macro}{gentryidx}
% This can be used inside a |gentry|; it expands to the first optional
% argument of the |gentry| environment, which is either blank or the
% requested fixed index of the entry.
% To get the actual shuffled index, use \cs{gentryidxu}.
% \end{macro}
%
% \begin{macro}{gentrycode}
% This can be used inside a |gentry|; it expands to the first mandatory
% argument of the |gentry| environment, which is the code for this
% entry (without the |gentry:| prefix).
% \end{macro}
%
% \cs{@@gentry} then reads in the optional title argument, storing it in the \cs{gentrytitle}
% macro to supply the unsorted index number and the current entry's code respectively.
%    \begin{macrocode}
\newcommand{\@@gentry}[1][]{%
  \def\gentrytitle{#1}%
  \stepcounter{gentryctr}%
%    \end{macrocode}
% For fixed entries, define a macro to hold the requested index
%    \begin{macrocode}
  \ifthenelse{\equal{\gentryidx}{}}{}{%
    \expandafter\xdef\csname fixedat\arabic{gentryctr}%
    \endcsname{\gentryidx}%
  }%
  \expandafter\global\expandafter\newtoks\expandafter{%
    \csname paratok\arabic{gentryctr}\endcsname}{ }%
  \Collect@Body\gentry@store%
%    \end{macrocode}
% This was supposed to discard any blank lines at the end of the
% |gentry| environment, but it still left odd spacing in the
% output somehow, and sometimes produced weird error messages about
% |\inaccessible| and |\head|. Removed for now.
%    \begin{macrocode}
%  \ignorespacesandallpars%
}
%    \end{macrocode}
% Store the collected environment contents for \cs{thegentries} to output:
% This uses the token register \cs{gentry$N$}, where $N$ is the
% unsorted index of this entry.
% Later, we'll change the values of the $N$ bit of \cs{gentry$N$}
% macros when we shuffle (the underlying token registers stay the same).
%
% This relies on \cs{refstepcounter}\marg{gentryctr} being expanded
% before this macro. The label |gentry:\gentrycode| is thus set to the
% current index of the final output entry.
%    \begin{macrocode}
\newcommand{\gentry@store}[1]{%
  \edef\head{\noexpand\begingroup\noexpand\gentryheader%
    {\arabic{gentryctr}}{\gentryidx}{\gentrycode}{\gentrytitle}%
    \noexpand\label{gentry:\gentrycode}%
  }%
  \global\expandafter\csname paratok\arabic{gentryctr}\endcsname=%
%    \end{macrocode}
% Output the header, the environment token list, flush any footnotes (if applicable) then the inter-gentry footer.    
%    \begin{macrocode}
  \expandafter{\head #1%
    \outputfootnotes@endgentry%
    \gentryfooter\endgroup}%
}%
%    \end{macrocode}
% \end{environment}
%
% \DescribeMacro{\gentry@footnotespergentry}
% This macro is executed inside the |gentry| group and sets up the
% commands to be run to allow footnotes. It is only expanded if
% \cs{if@gamebook@footnotes} is true, \textit{ie} the |footnotes|
% option is given.
%    \begin{macrocode}
\newcommand{\gentry@footnotespergentry}{}
%    \end{macrocode}
%
% \DescribeMacro{\thegentries}
%
% This macro shuffles the entries as required, then expands to them in
% the correct order.
%    \begin{macrocode}
\newcommand{\thegentries}{%
%    \end{macrocode}
% This command is in a group so that the output routine resets.
%
% The expansion of \cs{gentry@footnotespergentry} only happens if the
% |footnote| option was given; it sets up the output routine while the
% gamebook is running.
%    \begin{macrocode}
\begingroup%
  \if@gamebook@footnotes\gentry@footnotespergentry\fi%
%    \end{macrocode}
% To begin, let's record the number of entries. This may come in
% useful. Note that you can use this macro in the |gentry| environment,
% because that's not been expanded yet.
%    \begin{macrocode}
  \xdef\gentrycount{\arabic{gentryctr}}%
%    \end{macrocode}
%
% The next thing is to perform some surgery on LCG.
% This cuts out an \textbf{annoying} warning, hopefully more reliably
% than replacing the definition of \cs{p@stkeysr@nd}.
% The warning is simply that this package doesn't waste another
% counter every time it changes the random limits (which happens a lot
% during the Fisher-Yates shuffle):
%    \begin{macrocode}
  \WarningFilter{lcg}{Using an already existing counter rand}%
%    \end{macrocode}
%
% \subsection*{The output routine and end-of-page footnotes}
% The idea is to keep footnotes always on the same page as their mark where possible.
%
% \LaTeX\ does lots of fun things with the output routine, which we
% want to keep. So grab a copy of whatever the code is currently doing:
%
% Here I'm using the \cs{edef} trick to expand \cs{the}\cs{output} into a token register, because using a macro
% causes a weird error about an ``|{|'' after ``\cs{the}''.
%    \begin{macrocode}
\newtoks\gentry@oldoutput{}%
\edef\mytmp@{\noexpand\gentry@oldoutput={\the\output}}\mytmp@%
%    \end{macrocode}
% For the footnotes, if we reach the end of a page without outputting
% them, we need to flush them.
%
% \cs{output} is the output routine. It takes the page built up in
% \cs{box255}, annotates it with headers, footers etc, then ships it
% out. For our purposes, we only need to append the footnotes to the
% bottom of \cs{box255}.
%    \begin{macrocode}
\if@gamebook@footnotes\output={%
  \def\gentry@deferoutput{\the\gentry@oldoutput}%
%    \end{macrocode}
% The \cs{outputpenalty} tells us why the output routine was called; generally, it's invoked whenever
% a new floatable environment is generated, or when a page is full. Anything less than -1000 means that
% the page was filled, so we should add any footnotes only in this case.
%    \begin{macrocode}
  \ifnum\outputpenalty<-\@M\else%
    \if\gentryshouldoutput0%
      \unvbox255\def\gentry@deferoutput{}%
    \else%
      \expandafter\ifcsname footnotetoks\botmark\endcsname%
      \expandafter\if\expandafter\relax\expandafter%
      \detokenize\expandafter{\csname footnotetoks\botmark\endcsname}\relax\else%
      \global\setbox255=\vbox to \vsize{%
        \unvbox255\vfill\outputfootnotes@endpage}%
    \fi\fi%
  \fi\fi%
  \gentry@deferoutput%
%  \the\gentry@oldoutput
}\fi%
%    \end{macrocode}
%
%
% \subsection*{The Shuffling Algorithm}
%
% The basic shuffling algorithm is to first shuffle all entries,
% except for those marked with a fixed index, then to ge through the
% fixed-index entries in order and swap them into their final place.
%
% The original version of this package had a bug relating to multiple
% fixed-index entries (now fixed).
% In short, let $A$, $B$, and $C$ be indices; if $A<B$ and
% (unshuffled) entry number $A$ was fixed at (shuffled) 
% location $B$, while (unshuffled) entry number $B$ was fixed at
% (shuffled) location $C$, so during the ``shuffle,'' $A$ would be
% swapped with $B$, then $B$ would be swapped with $C$, resulting in
% $A$ appearing at $C$ in the text, not $B$ as requested. Because the
% unshuffled index doesn't appear in either the source or output
% document, this could be difficult to diagnose; the author simply saw
% one of their entries ending up in the wrong place.
%
% \TeX\ has well beyond 255 token registers these days, so don't
% bother to check that limit.
%
% The |LCG| package provides a suitable pseudo-random number generator.
% What we want is a repeatable series of disparate numbers, not an
% especially random one.
%
%\begin{enumerate}  
% \item Work out how many entries there are~($N$). Provided \cs{thegentries} is
% called at the end, this is just the value of |gentryctr|
% \item Declare a set of token registers named \cs{paratoks}$n$, where $n$ is
% each integer 1\dots$N$ inclusive. These will hold the contents of the entry.
% \item Declare a set of macros named \cs{paraidx}$n$, where $n$ is
% each integer 1\dots$N$ inclusive, each initialised to $n$. These will hold the number of the entry.
% \item Shuffle elements $\{2:N-1\}$, in that array. For $i=2$ through $N-2$
% \begin{enumerate}
% \item Let $R$ be a random number between $i$ and $N-1$ inclusive
% \item If $R\ne i$ then swap macros \cs{gentrytoks}$R$ and \cs{gentrytoks}$i$
% \item If $R\ne i$ then swap macros \cs{paraIdx}$R$ and \cs{gentryidx}$i$
% \end{enumerate}
% \item If a jukebox index sort is requested, perform an optimisation
% pass (see below)
% \item For $i=1:n$, output token reg $i$
%\end{enumerate}  
%
% Define macros \cs{csname}~|paraIdx|$n$\cs{endcsname} containing the arabic
% original gentry number to be put out on the $n$th output gentry.
% \begin{macro}{gentrycount}
%    \begin{macrocode}
  \xdef\gentrycount{\arabic{gentryctr}}%
%    \end{macrocode}
% \end{macro}  
% we can reuse |gentryctr|; we've finished this set of paras and
% kept the total count.
%    \begin{macrocode}
  \setcounter{gentryctr}{0}%
%    \end{macrocode}
% If there are fewer than 3 entries, don't try and shuffle.
%    \begin{macrocode}
  \ifthenelse{\gentrycount<3}{}{%SHUFFLE START
  \whiledo{\not{\value{gentryctr}>\gentrycount}}{%
    \edef\gentryidxu{\arabic{gentryctr}}%
    \expandafter\xdef\csname paraIdx\gentryidxu\endcsname{\gentryidxu}%
    \typeout{DEFINED paraIdx\gentryidxu}%
    \stepcounter{gentryctr}%
  }%
%    \end{macrocode}
% \begin{macro}{gentryidxu}
% \changes{v1.1}{2021/08/24}{Access shuffled index}
% The \cs{gentryidxu} macro can be used inside a |gentry| to obtain the
% current arabic shuffled index of the entry.
% \end{macro}  
% \begin{macro}{gentryidxs}
% \changes{v1.1}{2021/08/24}{Access original index}
% The \cs{gentryidxs} macro can be used inside a |gentry| to obtain the
% arabic original unshuffled index of the entry (the first |gentry| is
% 1, and this counter resets after each expansion of \cs{thegentries}).
% \end{macro}
%    \begin{macrocode}
  \if@gamebook@shuffle%
  \setcounter{rand}{\gentrycount}%
  \addtocounter{rand}{-1}\edef%
  \stoppoint{\arabic{rand}}%
%    \end{macrocode}
% First, shuffle everything that isn't fixed down.
% Don't renumber para~1 or \cs{gentrycount}; Fisher-Yates-shuffle the rest
% NB: we stop at \cs{gentrycount}$-2$, because \cs{gentrycount}$-1$ would only shuffle
% with itself.
%    \begin{macrocode}
  \setcounter{gentryctr}{2}%
  \chgrand[last=\stoppoint]%
  \whiledo{\value{gentryctr}<\stoppoint}{%
%    \end{macrocode}
% If this is to be swapped with a fixed position, skip it
%    \begin{macrocode}
    \edef\gentryidxu{\arabic{gentryctr}}%
    \expandafter\ifcsname fixedat\gentryidxu\endcsname%
      \gamebook@info{Not shuffling \gentryidxu; fixed pos}%
      \stepcounter{gentryctr}%
    \else%
      \gamebook@info{Shuffling \gentryidxu}%
      \stepcounter{gentryctr}%
      \edef\nextidx{\arabic{gentryctr}}%
%    \end{macrocode}
% Roll the dice. If we've hit an entry with fixed position, we must
% skip it, or it would end up being swapped out into
% |fixedat\arabic{rand}| instead.
%    \begin{macrocode}
      \chgrand[first=\nextidx]%
      \rand%
      \expandafter\ifcsname fixedat\arabic{rand}\endcsname\else%
        \gamebook@info{Shuffling \gentryidxu to \arabic{rand}}%
        \macroswap{paraIdx\gentryidxu}{paraIdx\arabic{rand}}%
      \fi%
    \fi%
  }%
%    \end{macrocode}
% Now move fixed entries into their final place:
%    \begin{macrocode}
  \setcounter{gentryctr}{2}%
  \whiledo{\not{\value{gentryctr}>\stoppoint}}{%
    \edef\gentryidxu{\arabic{gentryctr}}%
    \expandafter\ifcsname fixedat\gentryidxu\endcsname%
      \expandafter\edef\expandafter\mydest\expandafter%
        {\expandafter\csname fixedat\gentryidxu\endcsname}%
        \gamebook@info{MOVING FIXED GAMEBOOK ENTRY INTO PLACE: \gentryidxu -> \mydest}%
        \macroswap{paraIdx\gentryidxu}%
                {paraIdx\expandafter\csname fixedat\gentryidxu\endcsname}%
%    \end{macrocode}
% Edge case: It's possible that we also have |fixedat\mydest|; in which case that
% would be messed up with the reshuffling. So we need to rename that
% to |fixedat\gentryidxu| as well, then reprocess this index
%    \begin{macrocode}
      \expandafter\ifcsname fixedat\mydest\endcsname%
        \macroswap{fixedat\gentryidxu}{fixedat\mydest}%
        \expandafter\global\expandafter\let\csname fixedat\mydest\endcsname\@undefined%
        \addtocounter{gentryctr}{-1}%
      \fi%
%    \end{macrocode}
% if we are doing a jukebox shuffle, remember which final entries are fixed, so
% they don't get moved.
%    \begin{macrocode}
      \ifcsname gamebook@jukebox\endcsname%
        \expandafter\def\csname fixedto\mydest\endcsname{}%
      \fi%
    \fi%
    \stepcounter{gentryctr}%
  }%
%    \end{macrocode}
% The jukebox shuffle requires an extra pass. This must come after
% moving fixed entries into their final place, to allow us to compare
% the initial indicies.
% We make a reasonable effort:\begin{itemize}
% \item $t$ is the current index
% \item $u$ is the next index
% \item $r$ is a random index after $u$ (make 3 attempts to find a
% non-fixed $r$)
% \item if abs$(t-u)=1$ and $r$ is not fixed, then swap $u$ and $r$
% \textbf{if} $u$ is not fixed; otherwise:
% \item if abs$(t-u)=1$ and $r$ is not fixed, then swap $t$ and $r$
% \textbf{if} $t$ is not fixed;
% \item otherwise, give up.  
% \end{itemize}
%    \begin{macrocode}
  \ifcsname gamebook@jukebox\endcsname%
  \ifnum\gentrycount<6%
    \gamebook@info{Jukebox pass skipped; too few entries}%
  \else%
  \setcounter{gentryctr}{2}%
  \whiledo{\not{\value{gentryctr}=\stoppoint}}{%
    \edef\gentryidxt{\arabic{gentryctr}}%
    \edef\curpos{\csname paraIdx\gentryidxt\endcsname}%
    \stepcounter{gentryctr}%
    \edef\gentryidxu{\arabic{gentryctr}}%
    \edef\nextpos{\csname paraIdx\gentryidxu\endcsname}%
    \edef\pdiff{\the\numexpr\curpos+\nextpos}%
    \ifnum\pdiff<0\edef\pdiff{\the\numexpr-\pdiff}\fi%
    \edef\sdiff{\the\numexpr\curpos-\nextpos}%
    \ifnum\sdiff<0\edef\sdiff{\the\numexpr-\sdiff}\fi%
    \ifnum\pdiff<\sdiff\relax\def\thediff{\pdiff}\else\def\thediff{\sdiff}\fi%
    \ifnum\thediff=1%
      \gamebook@info{Jukebox: entries are too close: %
        \gentryidxt,\gentryidxu\space (original \curpos,\nextpos)}%
      \chgrand[first=\numexpr\gentryidxu+1]%
      \rand%
      \expandafter\ifcsname fixedto\arabic{rand}\endcsname\rand\fi%
      \expandafter\ifcsname fixedto\arabic{rand}\endcsname\rand\fi%
      \expandafter\ifcsname fixedto\arabic{rand}\endcsname%
        \gamebook@info{Can't reshuffle: failed to find non-fixed index}%
      \else%
        \ifcsname fixedto\gentryidxu\endcsname%
          \ifcsname fixedto\gentryidxt\endcsname%
          \gamebook@info{Can't reshuffle: \curpos\space and %
            \nextpos\space are both fixed.}%
          \else%
            \gamebook@info{Reshuffling \gentryidxt\space to \arabic{rand}}%
            \macroswap{paraIdx\gentryidxt}{paraIdx\arabic{rand}}%
          \fi%
        \else%
          \gamebook@info{Reshuffling \gentryidxu\space to%
            \arabic{rand} (alt)}%
          \macroswap{paraIdx\gentryidxu}{paraIdx\arabic{rand}}%
    \fi\fi\fi%
  }%
  \fi\fi%
  }\fi%SHUFFLE END
%    \end{macrocode}
% Now we can output the |gentry| token registers to let \LaTeX\ do its thing:
%    \begin{macrocode}
  \gamebook@info{Shuffled! Gentry order:}%
  \setcounter{gentryctr}{1}%
  \whiledo{\not{\value{gentryctr}>\gentrycount}}{%
    \edef\gentryidxu{\arabic{gentryctr}}%
    \gamebook@info{\gentryidxu -> \csname paraIdx\gentryidxu\endcsname}%
    \stepcounter{gentryctr}%
  }%
%    \end{macrocode}
%
% We can reuse the same counter again to output
%    \begin{macrocode}
  \setcounter{gentryctr}{0}%
  \gamebook@info{Outputting \gentrycount\space gamebook entries}%
  \whiledo{\value{gentryctr}<\gentrycount}{%
%    \end{macrocode}
% Use refstepcounter in the loop to allow \cs{label} to work as expected.
%    \begin{macrocode}
    \refstepcounter{gentryctr}%
%    \end{macrocode}
% Keep last entry on its own page
%    \begin{macrocode}
    \ifthenelse{\value{gentryctr}=\gentrycount}{%
      \gamebook@beforelast%
    }{}%
    \edef\gentryidxu{\arabic{gentryctr}}%
    \xdef\gentryidxs{\csname paraIdx\gentryidxu\endcsname}%
    \gamebook@info{Output gentry \gentryidxu\ of \gentrycount,%
      original idx \gentryidxs}%
%    \end{macrocode}
% Output the stored entry body, stripping any extraneous space:
%    \begin{macrocode}
    \the\csname paratok\gentryidxs\endcsname%
  }%
%    \end{macrocode}
% Finally, we clear the registers and reset the counter in case we want to start again
% (NB: fixedto is scope to the current block only, so no need to clear that)
%    \begin{macrocode}
  \gamebook@info{All gamebook entries added to main vertical list}%
  \setcounter{gentryctr}{1}%
  \whiledo{\not{\value{gentryctr}>\gentrycount}}{%
    \edef\gentryidxu{\arabic{gentryctr}}%
    \expandafter\global\expandafter\let%
      \csname paratok\gentryidxu\endcsname\@undefined%
    \expandafter\ifcsname fixedat\gentryidxu\endcsname%
      \expandafter\global\expandafter\let%
        \csname fixedat\gentryidxu\endcsname\@undefined%
    \fi%
    \stepcounter{gentryctr}%
  }%
  \setcounter{gentryctr}{0}%
  \eject%
  \endgroup%
}%
%    \end{macrocode}
% The final \cs{eject} ensures that the output routine has flushed as
% many pages as it can, before the output routine is reset again.
%
% \DescribeMacro{\gentryheader}\marg{counterIdx}\marg{fixedIdx}\marg{code}\marg{title}
%
% This macro is called before outputting the header. Its
% job is to format whatever header information the user wants to see
% on each entry; generally, this will be the page number.
%
% It takes 4 parameters:
%\begin{enumerate}
%\item The \textit{unshuffled} index value (to help an author find an
% entry in the original text)
%\item The \textit{fixed} index value, if any; otherwise an empty argument
%\item The user-entered unique code for this entry
%\item The user-supplied title, if any; otherwise an empty argument
%\end{enumerate}
%
% The value of the output index can be obtained with |\arabic{gentryctr}|
%    \begin{macrocode}
\newcommand{\gentryheader}[4]{%
  \noindent\textbf{\Huge\arabic{gentryctr}\large\ #4}%
  \nopagebreak%
  \vspace{0.3em}%
  \nopagebreak%
  \par%
  \marginpar{#3}%
}%
%    \end{macrocode}
%
% \begin{macro}{gentryshouldoutput}
% \changes{v1.3}{2022/05/26}{Suppress short pages option}
% This macro \textbf{may} be called from the output routine, and can
% be used to suppress page breaks. It was added because it proves
% fairly easy to write custom divider routines that can produce blank
% pages. If it expands to the number 1, then the page will be output;
% if it expands to 0, it will not be.
%
% For example, the following will prevent pages that are less than 80\% full.
%
% \verb!\renewcommand{\gentryshouldoutput}{%!
%
% \verb!   \ifdim\pagetotal>0.8\pagegoal\relax1\else 0\fi}!
%
% \textbf{Caution:} if this macro continually tests false, then
% material will eventually be discarded from the main vertical list to
% ensure that \TeX\ can complete the output of the document. If you
% redefine this macro, make sure to check out output carefully for
% missing text.
%    \begin{macrocode}
\newcommand{\gentryshouldoutput}{1}
%    \end{macrocode}
% \end{macro}
%
% \DescribeMacro{\gentryfooter}
% This takes no arguments and is simply expanded after the entry is typeset. The default adds some vertical
% space and a simple separator.
%
%    \begin{macrocode}
\newcommand{\gentryfooter}{%
  \par\vspace{2em}\centerline{---}\vspace{2em plus 1in}\par%
}%
%    \end{macrocode}
%
% \DescribeMacro{\ignorespacesandallpars@}
% This is a technique described on StackExchange\footnote{\small https://tex.stackexchange.com/questions/179016/\\\mbox{~}\hspace{1in}ignore-spaces-and-pars-after-an-environment\#179034}.
%
% This is used to ensure that extra space at the start and end of the |gentry|
% environment is ignored.
%
%    \begin{macrocode}
\def\ignorespacesandallpars@{%
  \@ifnextchar\par%
    {\expandafter\ignorespacesandallpars\@gobble}%
    {}%
}%
%    \end{macrocode}
%
% \DescribeMacro{\@footnotetext}\marg{token register}
% Token registers to hold each footnote.
%
% Each footnote is held in its own token register, so that we can control
% which footnotes appear on each page.
%
% \LaTeX\ would normally use saveboxes to store footnotes, but I
% prefer to hold off on expanding them until we know which page
% they're to be expanded on, which saves some difficulties (and
% perhaps creates others).
%
% To start with, modify \LaTeX's \cs{@makefnmark} to output a mark for
% the footnote, holding its index. This tells the output routine which
% footnotes to include at the end of the page. The
% \cs{in@out} macro is only defined when outputting footnotes, so
% suppresses this mark when we don't need it.
%    \begin{macrocode}
\g@addto@macro{\@makefnmark}{%
  \ifcsname in@out\endcsname\else%
    \mark{\arabic{\@mpfn}}%
  \fi%
}%
%    \end{macrocode}
%
% The \cs{footnotetext@save} is a convenience that takes the the
% footnote counter value in the first argument, the token register as
% a second argument, and the token list (footnote text) as the third.
%    \begin{macrocode}
\g@addto@macro{\gentry@footnotespergentry}{%
\newcommand{\@footnotetext@save}[3]{%
  \global\newtoks#2{}%
  \global#2={\noindent\@footnotemark{}#3}%
}}%
%    \end{macrocode}
% \cs{@footnotetext@save} unpacks the arguments for
% \cs{@footnotetext@save@}. This one uses the value of the \cs{@mpfn}
% counter to unpack the footnote counter, the token register (full csname) as
% the first argument, and the token list (footnote text) as the second.
%    \begin{macrocode}
\g@addto@macro{\gentry@footnotespergentry}{%
\newcommand{\@footnotetext@save@}[2]{%
  \edef\@tmp@{\expandafter\arabic{\@mpfn}}%
  \expandafter\@footnotetext@save\expandafter{\@tmp@}{#1}{#2}%
}}%
%    \end{macrocode}
% Footnotes are built in a group, so that \cs{in@out} can be defined
% locally to indicate that footnotes are being built (which stops
% spurious \cs{mark}s).
%    \begin{macrocode}
\g@addto@macro{\gentry@footnotespergentry}{%
\renewcommand{\@footnotetext}[1]{%
  \begingroup%
  \def\in@out{}% flag that we're building footnotes
  \edef\@tmp{\expandafter\csname footnotetoks\arabic{\@mpfn}\endcsname}%
  \expandafter\@footnotetext@save@\expandafter{\@tmp}{%
    % This fixes up the use of \cs{footnotemark} within footnotes:
    #1}%
  \endgroup%
}}%
%    \end{macrocode}
%
% \DescribeMacro{\outputfootnotes}\marg{maxIdx}
%
% Command to output the footnotes, which are just in \cs{footnotetoks$N$}
% for now.
%
% An end user can call this at any point in the text to set out the footnotes.
%
% This macro takes one argument, being the maximum index of footnotes to output, inclusive.
% Footnotes after this index will be excluded. If not provided, the
% value provided by the counter \cs{@mfpn}, which is the default
% \LaTeX\ counter, will be used.
%    \begin{macrocode}
\newcounter{fncounter}%
\newcommand{\outputfootnotes}[1]{%
  \begingroup%
  \def\in@out{}% flag that footnotes are outputting; suppresses marks
  \setcounter{fncounter}{1}%
%    \end{macrocode}
% Called in vertical mode, and don't want to throw a page break.
%
% I'm not sure how \LaTeX\ renders the footnote rule exactly, so I'm just using my own.
%    \begin{macrocode}
  \def\footnote@rule{%
%    \end{macrocode}
% The next line comes from the TUGboat suggestions, and protect against various user changes
%    \begin{macrocode}
    \leftskip=0pt\rightskip=0pt\interlinepenalty=1000%
    \penalty-1000%
    \vspace{1pt plus 2pt minus 0.5pt}%
    \hspace{-0.5in}\rule{1.5in}{0.4pt}\\%
  }%
%    \end{macrocode}
% Invoke the output routine. This attempts to stop the output routine
% being invoked while we're adding the footnote rule, which could
% cause a blank footnote to appear.
%
% Instead, it means some footnotes could appear on the subsequent
% page. To guard against this, footnotes are set into a \cs{vbox}
% to prevent page breaking.
%
% The macro \cs{footnote@rule} is output before each footnote. The
% first time, it outputs a divider rule; subsequently, it just throws
% a new paragraph to keep footnotes on separate lines.
%
% The next line comess from the TUGboat suggestions, and protect against various user changes
%    \begin{macrocode}
    \outputpenalty=-\@MM\break%
    \vbox{%
      \whiledo{\not{\value{fncounter}>#1}}{%
%    \end{macrocode}
% We are now looping over all possible entry numbers, in
% order. Some will already have been output, but we check them all
% anyway (it doesn't take much time).
%
% First, we check if the macro exists. If it does, it must contain the
% value of the token register.
%
% Next, we make sure that the token register has contents.
%
%    \begin{macrocode}
        \expandafter\ifcsname footnotetoks\arabic{fncounter}\endcsname%
        \edef\tmp@@{\csname footnotetoks\arabic{fncounter}\endcsname}%
%    \end{macrocode}
% This \cs{detokenize} black magic tests if a token reg is actually
% empty\footnote{\texttt{https://tex.stackexchange.com/questions/263733/\\\mbox{~}\hspace{1in}whats-the-best-practice-way-to-test-whether-parameter-is-empty}}
%    \begin{macrocode}
        \expandafter\if\expandafter\relax\expandafter%
          \detokenize\expandafter{\the\tmp@@}\relax\else%
        \footnote@rule\gdef\footnote@rule{\ifvmode\else\par\fi}%
        \interlinepenalty\interfootnotelinepenalty%
%    \end{macrocode}
% Actually output the footnotes.
%
% Having output the footnote, clear the token register (to save
% memory), then use \cs{let} to clear the definition of the
% macro. This ensures that we don't try to output the same footnote
% twice (at page end and entry end).
%    \begin{macrocode}
        \expandafter\the\tmp@@%
        \global\expandafter\tmp@@={}%
        \expandafter\let\tmp@@\@undefined%
      \fi\fi%
      \stepcounter{fncounter}%
      }%
  }%
  \endgroup%
}%
%    \end{macrocode}
% \LaTeX\ works by building up a little more than a page, then calling the output routine.
% The output routine then decides where to put the page end from the built-up material.
% If there is a footnote mark, it could come after the page end, so we can't rely on the
% fact that there's footnote register to determine if we should output footnotes or not.
% The edge case is: if the footnote mark is held back for the next page, the footnote text
% would appear on the footer of the page where being built when the mark was expanded, which
% is the page before the footnote.
%
% One fix for this is to have each footnote mark in the main body of
% the text output a \cs{mark} containing the footnote's counter value,
% and output the footnotes only to \cs{botmark}, which is the last
% \cs{mark} actually typeset on the page.
%
% (see https://www.tug.org/TUGboat/Articles/tb11-4/tb30salomon.pdf,
% page 598, to read around footnotes).
%
% \DescribeMacro{\outputfootnotes@endgentry}
% This macro is called at the end of each |gentry|. It switches into
% vertical mode (as ending |gentry| doesn't throw a \cs{break} or line
% end by itself), then we output all footnotes so far (by reading the
% footnote counter \cs{@mpfn}, as this is ``sequential'').
%    \begin{macrocode}
\newcommand{\outputfootnotes@endgentry}{%
  \if@gamebook@footnotes%
  \nopagebreak\ifhmode\\\fi% get into vertical mode
  \nopagebreak\outputfootnotes{\arabic{\@mpfn}}%
  \fi%
}%
%    \end{macrocode}
% \DescribeMacro{\noentryfoot}
% This method simply suppresses footnotes at the end of the entry
% within the current group. This forces footnotes within the current
% group to be printed at the bottom of the page.
%
% This is useful in the case where \LaTeX\ expands the footnote at the
% end of the entry, then decides to put the page split between the
% footnote mark and the text.
%    \begin{macrocode}
\newcommand{\noentryfoot}{\def\outputfootnotes@endgentry{}}
%    \end{macrocode}
% \DescribeMacro{\outputfootnotes@endpage}
% This is called by the output routine at the end of each page. If
% \cs{botmark} has a value, then we can output all footnotes up to
% that index.
%    \begin{macrocode}
\g@addto@macro{\gentry@footnotespergentry}{%
\newcommand{\outputfootnotes@endpage}{%
  \expandafter\if\expandafter\relax\expandafter%
  \detokenize\expandafter{\botmark}\relax\else%
  \outputfootnotes{\botmark}%
  \fi%
}}%
%    \end{macrocode}
% \Finale
\endinput
